% 由 nr_audit/make_numbers.py 生成，勿手改。字段映射见 nr_audit/results/nr_numbers_map.json。
% 含 \times 的宏须在数学模式中使用。
\newcommand{\nrICtOneLow}{18.48}  % 西安（η = 0.001N）仅隔离控制的启动时刻 t1，I0 为拟合值的 0.1 倍
\newcommand{\nrICtOneFit}{16.20}  % 西安（η = 0.001N）仅隔离控制的启动时刻 t1，I0 为拟合值的 1 倍
\newcommand{\nrICtOneHigh}{13.93}  % 西安（η = 0.001N）仅隔离控制的启动时刻 t1，I0 为拟合值的 10 倍
\newcommand{\nrICtOneShift}{2.28}  % I0 改为 0.1 或 10 倍时 t1 的最大变化（η = 0.001N，两种控制相同）
\newcommand{\nrICMaxRelSstar}{3.1\times10^{-9}}  % 六条轨迹中 S_star 相对拟合初值的最大相对变化（η = 0.001N）
\newcommand{\nrICMaxRelDt}{1.6\times10^{-9}}  % 六条轨迹中 dt 相对拟合初值的最大相对变化（η = 0.001N）
\newcommand{\nrICMaxRelJ}{4.3\times10^{-9}}  % 六条轨迹中 J 相对拟合初值的最大相对变化（η = 0.001N）
\newcommand{\nrICMaxRelTotal}{4.5\times10^{-9}}  % 六条轨迹中 total_infections_to_clearance 相对拟合初值的最大相对变化（η = 0.001N）
\newcommand{\nrICtabQLowTOne}{18.48}  % 西安（η = 0.001N）固定人口初值诊断：q_only，I0 为拟合值的 0.1 倍，t1
\newcommand{\nrICtabQLowDt}{170.6301}  % 西安（η = 0.001N）固定人口初值诊断：q_only，I0 为拟合值的 0.1 倍，dt
\newcommand{\nrICtabQLowJ}{82.119}  % 西安（η = 0.001N）固定人口初值诊断：q_only，I0 为拟合值的 0.1 倍，J
\newcommand{\nrICtabQLowTend}{391.59}  % 西安（η = 0.001N）固定人口初值诊断：q_only，I0 为拟合值的 0.1 倍，t_end
\newcommand{\nrICtabQFitTOne}{16.20}  % 西安（η = 0.001N）固定人口初值诊断：q_only，I0 为拟合值的 1 倍，t1
\newcommand{\nrICtabQFitDt}{170.6301}  % 西安（η = 0.001N）固定人口初值诊断：q_only，I0 为拟合值的 1 倍，dt
\newcommand{\nrICtabQFitJ}{82.119}  % 西安（η = 0.001N）固定人口初值诊断：q_only，I0 为拟合值的 1 倍，J
\newcommand{\nrICtabQFitTend}{389.32}  % 西安（η = 0.001N）固定人口初值诊断：q_only，I0 为拟合值的 1 倍，t_end
\newcommand{\nrICtabQHighTOne}{13.93}  % 西安（η = 0.001N）固定人口初值诊断：q_only，I0 为拟合值的 10 倍，t1
\newcommand{\nrICtabQHighDt}{170.6301}  % 西安（η = 0.001N）固定人口初值诊断：q_only，I0 为拟合值的 10 倍，dt
\newcommand{\nrICtabQHighJ}{82.119}  % 西安（η = 0.001N）固定人口初值诊断：q_only，I0 为拟合值的 10 倍，J
\newcommand{\nrICtabQHighTend}{387.04}  % 西安（η = 0.001N）固定人口初值诊断：q_only，I0 为拟合值的 10 倍，t_end
\newcommand{\nrICtabOLowTOne}{18.48}  % 西安（η = 0.001N）固定人口初值诊断：opt，I0 为拟合值的 0.1 倍，t1
\newcommand{\nrICtabOLowDt}{218.1027}  % 西安（η = 0.001N）固定人口初值诊断：opt，I0 为拟合值的 0.1 倍，dt
\newcommand{\nrICtabOLowJ}{76.624}  % 西安（η = 0.001N）固定人口初值诊断：opt，I0 为拟合值的 0.1 倍，J
\newcommand{\nrICtabOLowTend}{439.07}  % 西安（η = 0.001N）固定人口初值诊断：opt，I0 为拟合值的 0.1 倍，t_end
\newcommand{\nrICtabOFitTOne}{16.20}  % 西安（η = 0.001N）固定人口初值诊断：opt，I0 为拟合值的 1 倍，t1
\newcommand{\nrICtabOFitDt}{218.1027}  % 西安（η = 0.001N）固定人口初值诊断：opt，I0 为拟合值的 1 倍，dt
\newcommand{\nrICtabOFitJ}{76.624}  % 西安（η = 0.001N）固定人口初值诊断：opt，I0 为拟合值的 1 倍，J
\newcommand{\nrICtabOFitTend}{436.79}  % 西安（η = 0.001N）固定人口初值诊断：opt，I0 为拟合值的 1 倍，t_end
\newcommand{\nrICtabOHighTOne}{13.93}  % 西安（η = 0.001N）固定人口初值诊断：opt，I0 为拟合值的 10 倍，t1
\newcommand{\nrICtabOHighDt}{218.1027}  % 西安（η = 0.001N）固定人口初值诊断：opt，I0 为拟合值的 10 倍，dt
\newcommand{\nrICtabOHighJ}{76.624}  % 西安（η = 0.001N）固定人口初值诊断：opt，I0 为拟合值的 10 倍，J
\newcommand{\nrICtabOHighTend}{434.51}  % 西安（η = 0.001N）固定人口初值诊断：opt，I0 为拟合值的 10 倍，t_end
\newcommand{\nrFitMseCommNew}{47.33}  % 拟合初值下 community_new 的均方误差
\newcommand{\nrFitMseQuarNew}{130.38}  % 拟合初值下 quarantine_new 的均方误差
\newcommand{\nrFitMseCommCum}{387.22}  % 拟合初值下 community_cum 的均方误差
\newcommand{\nrFitMseQuarCum}{471.84}  % 拟合初值下 quarantine_cum 的均方误差
\newcommand{\nrFitObjective}{1036.76}  % 四项均方误差之和（拟合准则）
\newcommand{\nrWtBaseAToApLo}{3.1354}  % base：结构由 A 变为 A' 的 r 括区下端（宽度 < 1e-3，所检验网格上）
\newcommand{\nrWtBaseAToApHi}{3.1361}  % base：结构由 A 变为 A' 的 r 括区上端（宽度 < 1e-3，所检验网格上）
\newcommand{\nrWtBaseAToApApx}{3.14}  % base：结构由 A 变为 A' 的 r 括区中点取两位小数（正文用语“约”，不是括区端点）
\newcommand{\nrWtBaseApToBLo}{3.2050}  % base：结构由 A' 变为 B 的 r 括区下端（宽度 < 1e-3，所检验网格上）
\newcommand{\nrWtBaseApToBHi}{3.2057}  % base：结构由 A' 变为 B 的 r 括区上端（宽度 < 1e-3，所检验网格上）
\newcommand{\nrWtBaseApToBApx}{3.21}  % base：结构由 A' 变为 B 的 r 括区中点取两位小数（正文用语“约”，不是括区端点）
\newcommand{\nrWtBaseBToCLo}{13.1011}  % base：结构由 B 变为 C 的 r 括区下端（宽度 < 1e-3，所检验网格上）
\newcommand{\nrWtBaseBToCHi}{13.1018}  % base：结构由 B 变为 C 的 r 括区上端（宽度 < 1e-3，所检验网格上）
\newcommand{\nrWtBaseBToCApx}{13.10}  % base：结构由 B 变为 C 的 r 括区中点取两位小数（正文用语“约”，不是括区端点）
\newcommand{\nrWtXianAToApLo}{2.5477}  % xian：结构由 A 变为 A' 的 r 括区下端（宽度 < 1e-3，所检验网格上）
\newcommand{\nrWtXianAToApHi}{2.5482}  % xian：结构由 A 变为 A' 的 r 括区上端（宽度 < 1e-3，所检验网格上）
\newcommand{\nrWtXianAToApApx}{2.55}  % xian：结构由 A 变为 A' 的 r 括区中点取两位小数（正文用语“约”，不是括区端点）
\newcommand{\nrWtXianApToBLo}{2.7086}  % xian：结构由 A' 变为 B 的 r 括区下端（宽度 < 1e-3，所检验网格上）
\newcommand{\nrWtXianApToBHi}{2.7092}  % xian：结构由 A' 变为 B 的 r 括区上端（宽度 < 1e-3，所检验网格上）
\newcommand{\nrWtXianApToBApx}{2.71}  % xian：结构由 A' 变为 B 的 r 括区中点取两位小数（正文用语“约”，不是括区端点）
\newcommand{\nrWtXianBToCLo}{7.5108}  % xian（η = 0.001N）：结构由 B 变为 C 的 r 括区下端（宽度 < 1e-3）
\newcommand{\nrWtXianBToCHi}{7.5118}  % xian（η = 0.001N）：结构由 B 变为 C 的 r 括区上端（宽度 < 1e-3）
\newcommand{\nrWtXianBToCApx}{7.51}  % xian（η = 0.001N）：结构由 B 变为 C 的 r 括区中点取两位小数（正文用语“约”）
\newcommand{\nrConvTolMaxRel}{3\times10^{-12}}  % 积分容差三档间核心量的最大相对变化（阈值策略的全部轨迹）
\newcommand{\nrConvNzMaxRel}{1\times10^{-13}}  % N_z = 4001 与 8001 的 Δt、J 最大相对差
\newcommand{\nrConvNxMaxGap}{1\times10^{-13}}  % 控制网格上网格值低于根法值的最大相对量
\newcommand{\nrPlateauDevBase}{2.1\times10^{-11}}  % base 仅隔离控制开环完整方程阈值控制期的 max|I-η|（基准档）
\newcommand{\nrPlateauDevXian}{6.9\times10^{-11}}  % xian 仅隔离控制开环完整方程阈值控制期的 max|I-η|（基准档）
\newcommand{\nrConvNzAllMaxRel}{7.3\times10^{-13}}  % N_z = 1001、2001、4001 与 8001 的 Δt、J 最大相对差
\newcommand{\nrConvNzQuadMaxRel}{2.0\times10^{-13}}  % N_z = 8001 的 Simpson 与分段 quad 的 Δt、J 最大相对差
\newcommand{\nrConvPlateauMax}{4.6\times10^{-9}}  % 积分容差三档下阈值控制期 max|I-η| 的最大值
\newcommand{\nrConvExitMax}{1.8\times10^{-7}}  % 积分容差三档下退出时刻 S 与 S_c 之差的最大值
\newcommand{\nrConvMassMax}{1.3\times10^{-6}}  % 积分容差三档下人口守恒残差的最大值（含常规控制与 TDINN）
\newcommand{\nrTdinnTolChange}{1.5\times10^{-11}}  % 本模块收紧容差并把最大步长减半后，TDINN 总累计新增感染（清零时间）的变化
\newcommand{\nrIdentityMaxRel}{4\times10^{-14}}  % 累计新增感染恒等式残差/初始易感者总数的最大值
\newcommand{\nrRetLBase}{586.84}  % base：L = U(0) − S_c = S_0 + S_q(0) − S_c
\newcommand{\nrRetCinfBaseQ}{286.59}  % base：无回流（ω=0）q_only 的 t→∞ 累计新增感染 C_π(0)
\newcommand{\nrRetClrBaseQ}{285.7}  % base：无回流 q_only 计至清零时间的累计新增感染（表 2）
\newcommand{\nrRetCinfBaseA}{324.25}  % base：无回流（ω=0）alpha 的 t→∞ 累计新增感染 C_π(0)
\newcommand{\nrRetCinfBaseO}{307.50}  % base：无回流（ω=0）opt 的 t→∞ 累计新增感染 C_π(0)
\newcommand{\nrRetClrBaseO}{306.6}  % base：无回流 opt 计至清零时间的累计新增感染（表 2）
\newcommand{\nrRetCinfBaseC}{629.51}  % base：无回流（ω=0）c_only 的 t→∞ 累计新增感染 C_π(0)
\newcommand{\nrRetClrBaseC}{628.6}  % base：无回流 c_only 计至清零时间的累计新增感染（表 2）
\newcommand{\nrRetLXian}{10.1889}  % xian：L = U(0) − S_c = S_0 + S_q(0) − S_c
\newcommand{\nrRetCinfXianQ}{2.3779}  % xian：无回流（ω=0）q_only 的 t→∞ 累计新增感染 C_π(0)
\newcommand{\nrRetClrXianQ}{2.3779}  % xian：无回流 q_only 计至清零时间的累计新增感染（表 2）
\newcommand{\nrRetCinfXianA}{2.6390}  % xian：无回流（ω=0）alpha 的 t→∞ 累计新增感染 C_π(0)
\newcommand{\nrRetCinfXianO}{2.5348}  % xian：无回流（ω=0）opt 的 t→∞ 累计新增感染 C_π(0)
\newcommand{\nrRetClrXianO}{2.5348}  % xian：无回流 opt 计至清零时间的累计新增感染（表 2）
\newcommand{\nrRetCinfXianC}{3.8566}  % xian：无回流（ω=0）c_only 的 t→∞ 累计新增感染 C_π(0)
\newcommand{\nrRetClrXianC}{3.8566}  % xian：无回流 c_only 计至清零时间的累计新增感染（表 2）
\newcommand{\nrRetTermBaseSevenQ}{655.35}  % base：指数回流 1/ω=7 d，q_only，终止时（退出后首次 I<1）累计新增感染；来源见 source
\newcommand{\nrRetTermBaseSevenM}{657.51}  % base：指数回流 1/ω=7 d，opt_noreturn，终止时（退出后首次 I<1）累计新增感染；来源见 source
\newcommand{\nrRetTermBaseSevenC}{672.53}  % base：指数回流 1/ω=7 d，c_only，终止时（退出后首次 I<1）累计新增感染；来源见 source
\newcommand{\nrRetTermBaseFourteenQ}{619.48}  % base：指数回流 1/ω=14 d，q_only，终止时（退出后首次 I<1）累计新增感染；来源见 source
\newcommand{\nrRetTermBaseFourteenM}{620.50}  % base：指数回流 1/ω=14 d，opt_noreturn，终止时（退出后首次 I<1）累计新增感染；来源见 source
\newcommand{\nrRetTermBaseFourteenC}{665.80}  % base：指数回流 1/ω=14 d，c_only，终止时（退出后首次 I<1）累计新增感染；来源见 source
\newcommand{\nrRetTermBaseThirtyQ}{594.44}  % base：指数回流 1/ω=30 d，q_only，终止时（退出后首次 I<1）累计新增感染；来源见 source
\newcommand{\nrRetTermBaseThirtyM}{595.90}  % base：指数回流 1/ω=30 d，opt_noreturn，终止时（退出后首次 I<1）累计新增感染；来源见 source
\newcommand{\nrRetTermBaseThirtyC}{654.44}  % base：指数回流 1/ω=30 d，c_only，终止时（退出后首次 I<1）累计新增感染；来源见 source
\newcommand{\nrRetTermXianSevenQ}{10.5485}  % xian：指数回流 1/ω=7 d，q_only，终止时（退出后首次 I<1）累计新增感染；来源见 source
\newcommand{\nrRetTermXianSevenM}{10.5491}  % xian：指数回流 1/ω=7 d，opt_noreturn，终止时（退出后首次 I<1）累计新增感染；来源见 source
\newcommand{\nrRetTermXianSevenC}{10.5495}  % xian：指数回流 1/ω=7 d，c_only，终止时（退出后首次 I<1）累计新增感染；来源见 source
\newcommand{\nrRetTermXianFourteenQ}{10.4498}  % xian：指数回流 1/ω=14 d，q_only，终止时（退出后首次 I<1）累计新增感染；来源见 source
\newcommand{\nrRetTermXianFourteenM}{10.4520}  % xian：指数回流 1/ω=14 d，opt_noreturn，终止时（退出后首次 I<1）累计新增感染；来源见 source
\newcommand{\nrRetTermXianFourteenC}{10.4538}  % xian：指数回流 1/ω=14 d，c_only，终止时（退出后首次 I<1）累计新增感染；来源见 source
\newcommand{\nrRetTermXianThirtyQ}{10.3044}  % xian：指数回流 1/ω=30 d，q_only，终止时（退出后首次 I<1）累计新增感染；来源见 source
\newcommand{\nrRetTermXianThirtyM}{10.3031}  % xian：指数回流 1/ω=30 d，opt_noreturn，终止时（退出后首次 I<1）累计新增感染；来源见 source
\newcommand{\nrRetTermXianThirtyC}{10.3018}  % xian：指数回流 1/ω=30 d，c_only，终止时（退出后首次 I<1）累计新增感染；来源见 source
\newcommand{\nrRetTailBaseSeven}{1.581}  % base：平均停留7d三控制的最大剩余感染上界，显示向上取整
\newcommand{\nrRetTailBaseFourteen}{4.435}  % base：平均停留14d三控制的最大剩余感染上界，显示向上取整
\newcommand{\nrRetTailBaseThirty}{22.374}  % base：平均停留30d三控制的最大剩余感染上界，显示向上取整
\newcommand{\nrRetTailXianSeven}{10.738}  % xian：平均停留7d三控制的最大剩余感染上界，显示向上取整
\newcommand{\nrRetTailXianFourteen}{15.359}  % xian：平均停留14d三控制的最大剩余感染上界，显示向上取整
\newcommand{\nrRetTailXianThirty}{37.440}  % xian：平均停留30d三控制的最大剩余感染上界，显示向上取整
\newcommand{\nrRetSolverTolRel}{2.6\times10^{-11}}  % 固定来源54条容差对照中收紧后的核心量最大相对变化
\newcommand{\nrCostRcontBase}{3.1358}  % base：连续结构定理的权重界 r_cont(b) = 16/(b+sqrt(b^2-8b+8))^2
\newcommand{\nrCostRhatBase}{3.2051}  % base：前段仅隔离控制全局最优的权重界 r̂(b) = (S_sw/(S_sw-S_c))^2
\newcommand{\nrCostGapOneBase}{2.3}  % base：r = 1 时 100(1 - J_min/J_q-only)，即成本最优控制比仅隔离控制低的百分比
\newcommand{\nrCostRcontXian}{2.5477}  % xian：连续结构定理的权重界 r_cont(b) = 16/(b+sqrt(b^2-8b+8))^2
\newcommand{\nrCostRhatXian}{2.7090}  % xian：前段仅隔离控制全局最优的权重界 r̂(b) = (S_sw/(S_sw-S_c))^2
\newcommand{\nrCostGapOneXian}{3.0}  % xian（η = 0.001N）：r = 1 时 100(1 - J_min/J_q-only)
\newcommand{\nrRwBaseQDt}{19.21}  % base，1/ω = 14 d：仅隔离控制的控制时长
\newcommand{\nrRwBaseQJ}{2.823}  % base，1/ω = 14 d：仅隔离控制的成本 J（w_c = 1、w_q = 2）
\newcommand{\nrRwBaseQInf}{619.48}  % base，1/ω = 14 d：仅隔离控制计至 T_cut 的累计新增感染
\newcommand{\nrRwBaseODt}{20.69}  % base，1/ω = 14 d：回流下的成本最优控制的控制时长
\newcommand{\nrRwBaseOJ}{2.765}  % base，1/ω = 14 d：回流下的成本最优控制的成本 J（w_c = 1、w_q = 2）
\newcommand{\nrRwBaseOInf}{619.90}  % base，1/ω = 14 d：回流下的成本最优控制计至 T_cut 的累计新增感染
\newcommand{\nrRwBaseMDt}{21.81}  % base，1/ω = 14 d：无回流成本最优律的控制时长
\newcommand{\nrRwBaseMJ}{2.842}  % base，1/ω = 14 d：无回流成本最优律的成本 J（w_c = 1、w_q = 2）
\newcommand{\nrRwBaseMInf}{620.50}  % base，1/ω = 14 d：无回流成本最优律计至 T_cut 的累计新增感染
\newcommand{\nrRwBaseCDt}{38.47}  % base，1/ω = 14 d：仅减少接触的控制时长
\newcommand{\nrRwBaseCJ}{12.579}  % base，1/ω = 14 d：仅减少接触的成本 J（w_c = 1、w_q = 2）
\newcommand{\nrRwBaseCInf}{665.80}  % base，1/ω = 14 d：仅减少接触计至 T_cut 的累计新增感染
\newcommand{\nrRwBaseFrontPct}{11}  % base：回流下成本最优控制 λ 首次低于 0.99 的控制期比例（仅隔离前段）
\newcommand{\nrRwBaseMigGapPct}{2.8}  % base：沿用无回流成本最优律比回流下最低成本高出的百分比
\newcommand{\nrRwBaseL}{587}  % base：U(0) - S_c
\newcommand{\nrRwBaseDiffPct}{7.5}  % base，1/ω = 14 d：仅减少接触比仅隔离控制多出的累计感染（%）
\newcommand{\nrRwBaseDiffPctNoret}{120}  % base，无回流：仅减少接触比仅隔离控制多出的累计感染（%，计至清零）
\newcommand{\nrRwBaseQDtNoret}{5.90}  % base，无回流：仅隔离控制的控制时长
\newcommand{\nrRwBaseUMin}{0.55}  % base：各轨迹终止时 (S+S_q)/S_c 的最小值
\newcommand{\nrRwBaseUMax}{0.81}  % base：各轨迹终止时 (S+S_q)/S_c 的最大值
\newcommand{\nrRwXianQDt}{815.3}  % xian，1/ω = 14 d：仅隔离控制的控制时长
\newcommand{\nrRwXianQJ}{409.39}  % xian，1/ω = 14 d：仅隔离控制的成本 J（w_c = 1、w_q = 2）
\newcommand{\nrRwXianQInf}{10.4172}  % xian，1/ω = 14 d：仅隔离控制计至 T_cut 的累计新增感染
\newcommand{\nrRwXianODt}{1002.9}  % xian，1/ω = 14 d：回流下的成本最优控制的控制时长
\newcommand{\nrRwXianOJ}{369.24}  % xian，1/ω = 14 d：回流下的成本最优控制的成本 J（w_c = 1、w_q = 2）
\newcommand{\nrRwXianOInf}{10.4179}  % xian，1/ω = 14 d：回流下的成本最优控制计至 T_cut 的累计新增感染
\newcommand{\nrRwXianMDt}{930.2}  % xian，1/ω = 14 d：无回流成本最优律的控制时长
\newcommand{\nrRwXianMJ}{374.24}  % xian，1/ω = 14 d：无回流成本最优律的成本 J（w_c = 1、w_q = 2）
\newcommand{\nrRwXianMInf}{10.4179}  % xian，1/ω = 14 d：无回流成本最优律计至 T_cut 的累计新增感染
\newcommand{\nrRwXianCDt}{1744.5}  % xian，1/ω = 14 d：仅减少接触的控制时长
\newcommand{\nrRwXianCJ}{616.64}  % xian，1/ω = 14 d：仅减少接触的成本 J（w_c = 1、w_q = 2）
\newcommand{\nrRwXianCInf}{10.4184}  % xian，1/ω = 14 d：仅减少接触计至 T_cut 的累计新增感染
\newcommand{\nrRwXianFrontPct}{17}  % xian：回流下成本最优控制 λ 首次低于 0.99 的控制期比例（仅隔离前段）
\newcommand{\nrRwXianMigGapPct}{1.4}  % xian：沿用无回流成本最优律比回流下最低成本高出的百分比
\newcommand{\nrRwXianL}{10.19}  % xian：U(0) - S_c
\newcommand{\nrRwXianDiffPct}{0.01}  % xian，1/ω = 14 d：仅减少接触比仅隔离控制多出的累计感染（%）
\newcommand{\nrRwXianDiffPctNoret}{65}  % xian，无回流：仅减少接触比仅隔离控制多出的累计感染（%，计至清零）
\newcommand{\nrRwXianQDtNoret}{170.63}  % xian，无回流：仅隔离控制的控制时长
\newcommand{\nrRwXianUMin}{0.92}  % xian：各轨迹终止时 (S+S_q)/S_c 的最小值
\newcommand{\nrRwXianUMax}{0.92}  % xian：各轨迹终止时 (S+S_q)/S_c 的最大值
